\chapter{Datenputzen}
\section{Verteilung der Ansprecher und Time-over-Threshhold}

Zunächst werden bei den Daten der Langzeitmessung in Abbildung \ref{fig:driftzeitspektrum_draht_raw} die Driftzeit gegen Drahtnummer als 2D Histogramm aufgetragen und die Driftzeit gegen die Time-over-Threshhold (TOT) in Abbildung \ref{fig:dt_tot_raw} aufgetragen. Zur besseren Einheitlichkeit wurden für diese Diagramme bereits die noch nicht erläuterten neuen Drahtnummern aus Abschnitt \ref{sec:drahtkorrelation} genutzt.

In diesen Auftragungen erkennt man jeweils klare Strukturen;

\jafpp{driftzeitspektrum_draht_raw}{Driftzeiten pro Draht}{dt_tot_raw}{Driftzeit gegen TOT, jeweils ungefiltert}

In Abbildung \ref{fig:driftzeitspektrum_draht_raw} ist bis zu Driftzeiten von circa $\SI{250}{\nano\second}$ eine Zusammenhängende Verteilung mit einer Mehrheit der Treffer zu erkennen, darüber hinaus nur noch ein ungleichmäßiger Untergrund. In Abbildung \ref{fig:dt_tot_raw} ist eine Zusammenhängende Struktur mit TOT von $\SI{150}{\nano\second}$ bis $\SI{500}{\nano\second}$ und Driftzeiten von weniger als $\SI{300}{\nano\second}$ zu erkennen.
Weiterhin findet man bei TOT nahe null eine Häufung sowie kontinuierlich wachsendes Rauschen zu hohen Driftzeiten und kleinen TOT.
Die Relevante Struktur in Abbildung \ref{fig:dt_tot_raw} ist die zuerst beschriebene zusammenhängende mit Driftzeiten von weniger als $\SI{300}{\nano\second}$ und TOT zwischen $\SI{150}{\nano\second}$ und $\SI{500}{\nano\second}$. Aus der Struktur in Abbildung \ref{fig:driftzeitspektrum_draht_raw} erscheint es Logisch, dass Driftzeiten von sehr viel mehr als $\SI{300}{\nano\second}$ mehrheitlich Rauschen sind.
Daraus leiten sich die Filterbedingungen ab, dass ein Treffer mindestens eine TOT von $\SI{100}{\nano\second}$ mit gleichzeitig einer maximalen Driftzeit von $\SI{400}{\nano\second}$ haben darf\footnote{Es wurden $\SI{400}{\nano\second}$ gewählt, da die Kante der Struktur nicht ganz scharf ist um nicht unnötig echte Treffer zu ignorieren}. Trägt man beide Abbildungen mit diesen Filtern erneut auf, erhält man die Abbildungen \ref{fig:driftzeitspektrum_draht} und \ref{fig:dt_tot}. Hier erkennt man die bereits beschriebenen Strukturen nun deutlicher, da weniger Rauschen vorhanden ist.

\jafpp{driftzeitspektrum_draht}{Driftzeiten pro Draht}{dt_tot}{Driftzeit gegen TOT, jeweils gefiltert}

\section{Drahtkorrelation}
\label{sec:drahtkorrelation}
Um eine korrekte Sortierung der Drähte zu gewährleisten wird in einem 2D Histogram aufgetragen, welche Drähte gemeinsam ansprechen. Zu sehen ist die in \ref{fig:drahtkorrelation_raw}. Dort erkennt man, dass die Maxima jeder Spalte abwechselnd nach um zwei Bins links und rechts von der Diagonalen entfernt liegen. Dies Spricht dafür, dass die Nummerierung der beiden Lagen vertauscht werden muss, damit bei steigender Drahtnummer eine gleiche Positionsänderung vorliegt. Verdeutlich haben wir dies - Skizzenhaft - in den Abbildungen \ref{fig:hexagons_pre_final} und \ref{fig:hexagons_correct_final}. Nach der Umnummerierung liegen die Maxima auf den Nebendiagonalen, was bedeutet, dass hauptsächlich aufeinander folgende Drahtnummern gemeinsam ansprechen. Da hauptsächlich nebeneinanderliegende Zellen ansprechen, zeigt dies die korrektheit der neuen Nummerierung.


\jafpp{drahtkorrelation_raw}{Drahtkorrelation vor}{drahtkorrelation_sortiert}{nach Umnummerierung der Drähte.}
\jafpp{hexagons_pre_final}{Drahtnummern, wie sie von der Software angegeben werden}{hexagons_correct_final}{nach der Umnummerierung.}
